\documentclass[11pt]{article}

\usepackage{amsmath,amssymb,amsthm,bm,mathtools}
\usepackage{geometry}
\geometry{margin=1in}
\usepackage{setspace}
\usepackage{enumitem}
\usepackage{booktabs}
\usepackage{hyperref}
\usepackage{float}

\onehalfspacing

\newcommand{\R}{\mathbb R}
\newcommand{\ind}{\mathbb I}
\newcommand{\norm}[1]{\left\lVert #1 \right\rVert}
\newcommand{\trans}{^{\mathsf T}}
\DeclareMathOperator{\logit}{logit}
\DeclareMathOperator{\Cat}{Categorical}
\DeclareMathOperator{\Dir}{Dirichlet}
\DeclareMathOperator{\IG}{IG}
\DeclareMathOperator{\IW}{IW}
\DeclareMathOperator{\tr}{tr}

\title{Joint Multilayered LSIRM with Latent-Position Skew-$t$ Mixture Clustering\\
Residual-Augmented NIW Predictive Allocation and Jain--Neal Split--Merge Moves}
\date{Model Description (v12\_skewt)}

\begin{document}
\maketitle

\section{Overview}

This document defines a joint multilayered Latent Space Item Response Model
(LSIRM) in which item clustering is performed directly on item latent
positions, and the within-cluster component distribution is a multivariate
\emph{skew-$t$} in the Sahu--Dey--Branco (SDB) parameterization.  Each item
latent position is assigned a finite skew-$t$ mixture prior,
\[
    b_j^{(\ell(j))}\mid c_j=l,\mu_l,\Sigma_l,\Delta_l
    \;\sim\; \mathrm{SkewT}_{\nu_t}(\mu_l,\Sigma_l,\Delta_l),
    \qquad d=2,
\]
where $c_j$ is the item cluster label, $(\mu_l,\Sigma_l)$ are the
cluster location and scatter, $\Delta_l\in\R^d$ is the per-cluster
\emph{skew direction}, and $\nu_t$ is a fixed Student-$t$ degrees-of-freedom
hyperparameter.  Setting $\Delta_l=0$ recovers an elliptical Student-$t$
component; letting $\nu_t\to\infty$ further recovers a Gaussian component.
The mixture prior is imposed only on the item latent positions; no
cluster-specific prior is placed on respondent ability, item difficulty, or
layer-specific nuisance parameters.

The model is fully joint: item cluster labels affect the posterior of item
positions through the mixture prior, and item positions affect the response
likelihood through the LSIRM distance term.  Two latent variables per item
make SDB skew-$t$ tractable as a Gaussian location--scale mixture:
\begin{itemize}[leftmargin=1.5em]
    \item a positive scale latent $w_q\sim\mathrm{Gamma}(\nu_t/2,\nu_t/2)$,
    inherited from the elliptical Student-$t$ representation
    of Lange, Little, and Taylor (1989);
    \item a positive location latent $U_q\mid w_q\sim N_+(0,1/w_q)$,
    a half-normal that drives the asymmetric mean shift of SDB skew-$t$.
\end{itemize}
Conditional on $(U_q,w_q)$, each $b_q$ is exactly Gaussian, so all
conjugate Normal--Inverse-Wishart (NIW) machinery for $(\mu_l,\Sigma_l)$
carries over verbatim once $b_q$ is replaced by the residual
$r_q=b_q-\Delta_{c_q}U_q$ in the weighted sufficient statistics.  The skew
direction $\Delta_l$ is updated by an additional multivariate Gaussian
regression Gibbs step.

The motivation for the skew extension is that LSIRM latent positions in
layers with asymmetric item difficulty distributions (e.g.\ count layers
with right-skewed dispersion, continuous layers with bounded support, ordinal
layers with non-uniform thresholds) cannot be captured by elliptical
components without inflating the cluster covariance and dragging the
cluster identification.  Allowing $\Delta_l\neq 0$ lets one cluster carry
both a center and a direction of asymmetric spread without splitting into
two sub-clusters.

\section{Indexing, dimensions, and observed data}

\begin{center}
\begin{tabular}{lll}
\toprule
Symbol & Meaning & Dimension/range \\
\midrule
$\ell$ & layer index & $\ell=1,\dots,L$ \\
$i$ & respondent index & $i=1,\dots,n$ \\
$j$ & local item index within layer $\ell$ & $j=1,\dots,P_\ell$ \\
$q$ & global item index across all layers & $q=1,\dots,P$, $P=\sum_\ell P_\ell$ \\
$\ell(q)$ & layer containing global item $q$ & $\{1,\dots,L\}$ \\
$j(q)$ & local item index of global item $q$ & $\{1,\dots,P_{\ell(q)}\}$ \\
$d$ & latent space dimension & $d=2$ in implementation \\
$K^\star$ & truncation level (maximum mixture components) & positive integer \\
$K_+$ & number of occupied components & $K_+=|\{l:n_l>0\}|$ \\
\bottomrule
\end{tabular}
\end{center}

\paragraph{Common-respondent assumption.}
The typical multilayered setting is assumed: every layer shares the same
$n$ respondents, and the respondent latent position $a_i\in\R^d$ is shared
across layers.  The layered notation $a_i^{(\ell)}$ is therefore understood
as a single $a_i$.  The only identifiability ambiguity that remains in the
latent space is a single global rigid motion of the whole configuration
$(a,b^{(1)},\dots,b^{(L)})$ together with $(\mu_l,\Sigma_l,\Delta_l)$.

The observed data are
\[
    Y^{(\ell)}=\{y_{ij}^{(\ell)}:(i,j)\in\mathcal O_\ell\},
    \qquad \ell=1,\dots,L,
\]
with $\mathcal O_\ell$ the set of observed respondent--item pairs in layer
$\ell$.

\section{LSIRM likelihood}

For respondent $i$ and item $j$ in layer $\ell$, the model uses the
following parameters:
\begin{itemize}[leftmargin=1.5em]
    \item $a_i\in\R^d$: shared respondent latent position.
    \item $b_j^{(\ell)}\in\R^d$: item latent position in layer $\ell$.
    \item $\alpha_i^{(\ell)}\in\R$: respondent effect in layer $\ell$.
    \item $\beta_j^{(\ell)}\in\R$: item difficulty for binary, continuous,
    and count layers.  For ordinal layers $\beta_j^{(\ell)}$ is replaced by
    an item-specific threshold vector
    $\beta_{j,1}^{(\ell)}<\cdots<\beta_{j,K_\ell-1}^{(\ell)}$ with
    $K_\ell$ ordered categories.
    \item $\gamma_\ell>0$: layer-specific distance weight.
    \item $\sigma_0^2>0$: continuous-layer error scale.
    \item $\nu_2>0$: continuous-layer Student-$t$ degrees of freedom.
    \item $\kappa_j>0$: item-specific count-layer overdispersion (negative
    binomial), one per count item $j$.
\end{itemize}
Let $\mathcal L_{\mathrm{ord}}\subseteq\{1,\dots,L\}$ index the ordinal
layers.

The layer-specific LSIRM linear predictor is
\begin{equation}
    \eta_{ij}^{(\ell)}
    =
    \alpha_i^{(\ell)}-\beta_j^{(\ell)}-\gamma_\ell\norm{a_i-b_j^{(\ell)}}_2,
    \qquad \ell\notin\mathcal L_{\mathrm{ord}},
    \label{eq:lsirm-eta}
\end{equation}
with the convention that for ordinal layers item difficulty enters via the
threshold vector, not the linear predictor:
\begin{equation}
    \eta_{ij}^{(\ell)}
    =
    \alpha_i^{(\ell)}-\gamma_\ell\norm{a_i-b_j^{(\ell)}}_2,
    \qquad \ell\in\mathcal L_{\mathrm{ord}}.
    \label{eq:lsirm-eta-ord}
\end{equation}

\subsection{Per-layer observation models}

\paragraph{Binary (logistic regression).}
For a binary layer $\ell$,
\begin{equation}
    y_{ij}^{(\ell)}\mid \eta_{ij}^{(\ell)}\sim\mathrm{Bernoulli}(p_{ij}^{(\ell)}),
    \qquad
    \logit p_{ij}^{(\ell)}=\eta_{ij}^{(\ell)}.
    \label{eq:lik-binary}
\end{equation}

\paragraph{Continuous (robust Student-$t$ via normal--gamma scale mixture).}
For a continuous layer $\ell$,
\begin{align}
    y_{ij}^{(\ell)}\mid \eta_{ij}^{(\ell)},\lambda_{ij}^{(\ell)}
    &\sim N\!\left(\eta_{ij}^{(\ell)},\;\sigma_0^2/\lambda_{ij}^{(\ell)}\right),
    \notag\\
    \lambda_{ij}^{(\ell)} &\sim\mathrm{Gamma}(\nu_2/2,\nu_2/2).
    \label{eq:lik-continuous}
\end{align}
Marginalizing over $\lambda_{ij}^{(\ell)}$ yields
$y_{ij}^{(\ell)}\mid\eta_{ij}^{(\ell)}\sim t_{\nu_2}(\eta_{ij}^{(\ell)},\sigma_0^2)$.

\paragraph{Count (negative binomial regression with item-specific overdispersion).}
For a count layer $\ell$ with item-specific overdispersion $\kappa_j>0$,
\begin{equation}
    y_{ij}^{(\ell)}\mid\eta_{ij}^{(\ell)},\kappa_j
    \sim\mathrm{NB}\!\left(r=1/\kappa_j,\,\mu=\exp\eta_{ij}^{(\ell)}\right),
    \label{eq:lik-count}
\end{equation}
parameterized so that $E[y]=\mu$ and $\mathrm{Var}[y]=\mu+\kappa_j\mu^2$.
The Poisson limit is $\kappa_j\to 0$.  Each $\kappa_j$ has a log-normal
prior:
\begin{equation}
    \log\kappa_j\sim N(\mu_{\log\kappa},\sigma_{\log\kappa}^2),
    \qquad j=1,\dots,P_\ell.
    \label{eq:kappa-prior}
\end{equation}

\paragraph{Ordinal (latent-space graded response model, LSGRM).}
For an ordinal layer $\ell\in\mathcal L_{\mathrm{ord}}$ with $K_\ell$
ordered categories and item thresholds
$\beta_{j,1}^{(\ell)}<\cdots<\beta_{j,K_\ell-1}^{(\ell)}$,
\begin{equation}
    \logit
    P(y_{ij}^{(\ell)}\geq k+1\mid\eta_{ij}^{(\ell)},\beta_{j,k}^{(\ell)})
    =\eta_{ij}^{(\ell)}-\beta_{j,k}^{(\ell)},
    \qquad k=1,\dots,K_\ell-1,
    \label{eq:lik-ord-cum}
\end{equation}
with $P(y\geq 1)=1$ and $P(y\geq K_\ell+1)=0$.  Category probabilities are
obtained by differencing:
\begin{equation}
    P(y_{ij}^{(\ell)}=k\mid\cdot)=P(y\geq k\mid\cdot)-P(y\geq k+1\mid\cdot).
    \label{eq:lik-ord}
\end{equation}

\subsection{Compact notation}

Each per-layer likelihood is summarized as
\begin{equation}
    y_{ij}^{(\ell)}\mid\eta_{ij}^{(\ell)},\xi_\ell
    \sim f_\ell\{y_{ij}^{(\ell)}\mid\eta_{ij}^{(\ell)},\xi_\ell\},
    \label{eq:lik-generic}
\end{equation}
with $\xi_\ell$ collecting all layer-specific nuisance parameters
($\sigma_0^2,\lambda_{ij}^{(\ell)}$ in continuous layers; $\{\kappa_j\}$ in
count layers; threshold vectors in ordinal layers; $\gamma_\ell$ in all
layers).  The full LSIRM log-likelihood is
\begin{equation}
    \ell_{\mathrm{LSIRM}}=
    \sum_{\ell=1}^L\sum_{(i,j)\in\mathcal O_\ell}
    \log f_\ell\{y_{ij}^{(\ell)}\mid\eta_{ij}^{(\ell)},\xi_\ell\}.
    \label{eq:lsirm-loglik}
\end{equation}

\section{Latent-position skew-$t$ mixture prior for items}
\label{sec:skewt-mixture}

\subsection{Global item position notation}

The global item latent position is
\begin{equation}
    z_q=b_{j(q)}^{(\ell(q))}\in\R^d,
    \qquad q=1,\dots,P.
\end{equation}
The clustering model is placed directly on $z_q$.

\subsection{SDB skew-$t$ hierarchical representation}

For $q=1,\dots,P$ and $l=1,\dots,K^\star$, the skew-$t$ mixture component is
expressed as a Gaussian location--scale mixture in two stages.  First the
scale latent and the location latent are introduced:
\begin{align}
    w_q
    &\sim\mathrm{Gamma}\!\bigl(\nu_t/2,\nu_t/2\bigr),
    \label{eq:w-prior}\\
    U_q\mid w_q
    &\sim N_+\!\bigl(0,1/w_q\bigr),
    \label{eq:U-prior}
\end{align}
where $N_+(0,1/w_q)$ denotes the half-normal distribution with density
$\sqrt{2w_q/\pi}\exp(-w_q U^2/2)$ on $U>0$.  Both latents are
\emph{item-specific}, indexed by the global item index $q$.

Conditional on the cluster label, the scale, and the location latent, the
item position is Gaussian:
\begin{equation}
    z_q\mid c_q=l,w_q,U_q,\mu_l,\Sigma_l,\Delta_l
    \sim
    N_d\!\bigl(\mu_l+\Delta_l U_q,\;\Sigma_l/w_q\bigr).
    \label{eq:z-given-latents}
\end{equation}
Marginalizing $U_q$ first gives the skew-Normal conditional on $w_q$;
marginalizing also $w_q$ gives the SDB skew-$t$ marginal,
\begin{equation}
    z_q\mid c_q=l,\mu_l,\Sigma_l,\Delta_l
    \sim
    \mathrm{SkewT}_{\nu_t}\!\bigl(\mu_l,\Sigma_l,\Delta_l\bigr).
    \label{eq:skewt-marginal}
\end{equation}
The cluster label and the mixture weights follow the standard finite-mixture
prior with a symmetric Dirichlet hyperprior:
\begin{align}
    c_q\mid\rho &\sim\Cat(\rho_1,\dots,\rho_{K^\star}),
    \label{eq:c-prior}\\
    \rho &\sim\Dir(e_0,\dots,e_0).
    \label{eq:rho-prior}
\end{align}

\paragraph{Limiting cases.}
The representation \eqref{eq:w-prior}--\eqref{eq:z-given-latents} contains
two ``recovery'' limits that are useful diagnostics:
\begin{itemize}[leftmargin=1.5em]
    \item $\Delta_l\to 0$:  the location latent $U_q$ drops out of the
    conditional mean, and \eqref{eq:skewt-marginal} reduces to the
    elliptical Student-$t$ distribution $t_{\nu_t}(\mu_l,\Sigma_l)$.
    \item $\nu_t\to\infty$:  the scale latent $w_q\to 1$ in probability and
    \eqref{eq:skewt-marginal} reduces to the skew-Normal distribution of
    Azzalini and Dalla Valle (1996), or further to $N_d(\mu_l,\Sigma_l)$ if
    also $\Delta_l\to 0$.
\end{itemize}

\subsection{Cluster-level priors}

The cluster location and scatter follow a Normal--Inverse-Wishart (NIW)
prior:
\begin{align}
    \Sigma_l
    &\sim\IW(\nu_0,S_0),
    \label{eq:Sigma-prior}\\
    \mu_l\mid\Sigma_l
    &\sim N_d\!\left(m_0,\;\frac{1}{\kappa_0}\Sigma_l\right),
    \qquad l=1,\dots,K^\star,
    \label{eq:mu-prior}
\end{align}
with hyperparameters $m_0\in\R^d$, $\kappa_0>0$, $\nu_0>d-1$, and
$S_0\in\R^{d\times d}$ positive definite.

The skew direction has an independent zero-mean isotropic Gaussian prior:
\begin{equation}
    \Delta_l\sim N_d\!\bigl(0,\,\tau_\Delta^2\,I_d\bigr),
    \qquad l=1,\dots,K^\star.
    \label{eq:Delta-prior}
\end{equation}
The prior scale $\tau_\Delta>0$ shrinks $\Delta_l$ toward zero; in the
limit $\tau_\Delta\to 0$ the model reduces to the elliptical
Student-$t$ mixture, so the prior \eqref{eq:Delta-prior} provides a natural
nested-model regularization.

\paragraph{Data-driven cluster scale: Wishart hyperprior on $S_0$.}
The NIW scale matrix $S_0$ is itself given a Wishart hyperprior so that the
cluster-covariance scale is learned from data rather than tuned:
\begin{equation}
    S_0\sim W(\nu_{S_0},\Lambda_{S_0}),
    \qquad \nu_{S_0}>d-1,\;\Lambda_{S_0}\succ 0.
    \label{eq:S0-hyperprior}
\end{equation}
With $\Sigma_l\mid S_0\sim\IW(\nu_0,S_0)$ for $l=1,\dots,K^\star$ this gives
a conjugate Wishart full conditional,
\begin{equation}
    S_0\mid\{\Sigma_l\}
    \sim
    W\!\left(
    \nu_{S_0}+K^\star\nu_0,\;
    \bigl(\Lambda_{S_0}^{-1}+\sum_{l=1}^{K^\star}\Sigma_l^{-1}\bigr)^{-1}
    \right),
    \label{eq:S0-update}
\end{equation}
sampled once per MCMC sweep after the $(\mu_l,\Sigma_l)$ draws.  Setting
$\nu_{S_0}\to\infty$ holds $S_0$ fixed at $E[S_0]$.

\paragraph{Modeling restriction.}
The cluster label $c_q$ enters only the prior for $z_q=b_{j(q)}^{(\ell(q))}$.
There is no cluster-specific prior on $\alpha_i^{(\ell)}$,
$\beta_j^{(\ell)}$, $\xi_\ell$, or any other parameter.  Cluster labels
therefore represent latent item geometry only.

\section{Priors for non-clustering LSIRM parameters}

\paragraph{Respondent latent position (fixed unit variance).}
For identifiability of the latent-space scale,
\begin{equation}
    a_i\sim N_d(0,I_d),\qquad i=1,\dots,n.
    \label{eq:a-prior}
\end{equation}
The unit variance is not sampled: the latent space has no intrinsic scale
because $a_i$ and $b_j^{(\ell)}$ can both be rescaled by an arbitrary
constant absorbed into $\gamma_\ell$.  Anchoring the respondent prior
variance at $1$ leaves $\gamma_\ell$ as the free scale parameter to be
estimated from data.

\paragraph{Respondent effect.}
Each layer's respondent effect has its own variance hyperparameter:
\begin{align}
    \alpha_i^{(\ell)}\mid\sigma_{\alpha,\ell}^2
    &\sim N(0,\sigma_{\alpha,\ell}^2),
    \qquad i=1,\dots,n,\,\ell=1,\dots,L,
    \label{eq:alpha-prior}\\
    \sigma_{\alpha,\ell}^2
    &\sim\IG(A_{\alpha,\ell},B_{\alpha,\ell}),
    \label{eq:sigma-alpha-prior}
\end{align}
allowing the respondent-effect scale to differ across layers (e.g.\ binary
versus count layers).

\paragraph{Item difficulty (fixed-scale prior).}
For non-ordinal layers,
\begin{equation}
    \beta_j^{(\ell)}\sim N(0,\sigma_{\beta,\ell}^2),
    \qquad j=1,\dots,P_\ell,\,\ell\notin\mathcal L_{\mathrm{ord}},
    \label{eq:beta-prior}
\end{equation}
with $\sigma_{\beta,\ell}^2$ a fixed hyperparameter (not sampled).  No
across-item shrinkage is imposed; each item's difficulty is recovered
directly from its own data.  The implementation uses
$\sigma_{\beta,\ell}=2$ (variance $4$).

\paragraph{Ordinal thresholds.}
For ordinal layers,
\begin{equation}
    \beta_{j,k}^{(\ell)}\sim N(0,\sigma_{\beta,\ell}^2),
    \qquad
    j=1,\dots,P_\ell,\,k=1,\dots,K_\ell-1,
    \label{eq:beta-thr-prior}
\end{equation}
restricted to $\beta_{j,1}^{(\ell)}<\cdots<\beta_{j,K_\ell-1}^{(\ell)}$.
The scale $\sigma_{\beta,\ell}^2$ is fixed; the implementation uses
$\sigma_{\beta,\ell}=2$.

\paragraph{Layer-specific nuisance parameters.}
The distance weight $\gamma_\ell>0$ in all layers is parameterized through
$\log\gamma_\ell\sim N(\mu_{\log\gamma_\ell},\sigma_{\log\gamma_\ell}^2)$
with fixed hyperparameters.  The continuous-layer error scale
$\sigma_0^2$ has an inverse-gamma prior $\IG(a_{\sigma_0},b_{\sigma_0})$;
$\nu_2$ is fixed.  The count-layer overdispersions $\{\kappa_j\}$ follow
the log-normal prior \eqref{eq:kappa-prior}.

The item latent position $b_j^{(\ell)}$ receives no additional independent
Gaussian prior beyond the skew-$t$ mixture prior in
\eqref{eq:z-given-latents}--\eqref{eq:rho-prior},
\eqref{eq:Sigma-prior}--\eqref{eq:Delta-prior}.

\section{Rotation and reflection invariance}
\label{sec:rotation-invariance}

Under the common-respondent assumption, the entire configuration
$(a,b^{(1)},\dots,b^{(L)})$ shares one coordinate system.  The only
identifiability ambiguity in the latent space is a single global orthogonal
transformation $Q\in O(d)$ together with a global translation $t\in\R^d$:
\begin{equation}
    a_i\mapsto Qa_i+t,
    \quad
    b_j^{(\ell)}\mapsto Qb_j^{(\ell)}+t,
    \quad
    \mu_l\mapsto Q\mu_l+t,
    \quad
    \Sigma_l\mapsto Q\Sigma_l Q\trans,
    \quad
    \Delta_l\mapsto Q\Delta_l.
    \label{eq:rigid-motion}
\end{equation}
The skew direction transforms as a free vector (rotation only, no
translation) because it captures the direction of within-cluster spread.
The half-normal latent $U_q$ is rotation-invariant (it is a scalar).

The LSIRM linear predictor in \eqref{eq:lsirm-eta}--\eqref{eq:lsirm-eta-ord}
depends on positions only through the Euclidean norm
$\norm{a_i-b_j^{(\ell)}}_2$, which is invariant under
\eqref{eq:rigid-motion}.  With $m_0=0$ and $S_0=s_0 I_d$, the NIW prior on
$(\mu_l,\Sigma_l)$ and the isotropic Gaussian prior on $\Delta_l$ in
\eqref{eq:Delta-prior} are also invariant in distribution under
$Q\in O(d)$.  Therefore the joint posterior on the clustering targets
($c$, $K_+$, the cluster sizes, the posterior similarity matrix
$C_{qr}=P(c_q=c_r\mid Y)$) is invariant under \eqref{eq:rigid-motion}.

\paragraph{Consequence for the sampler.}
Within-MCMC orientation alignment is not required for clustering
correctness.  A post-hoc Procrustes alignment is recommended only for
visualization or comparison with an external target configuration.  Such an
alignment is a deterministic post-processing step applied to stored
posterior draws; $\Delta_l$ is co-rotated as in \eqref{eq:rigid-motion}.

\section{Full joint posterior}

Collect the parameters as
\begin{align*}
    \Theta_{\mathrm{LSIRM}}
    &=\{a_i,\alpha_i^{(\ell)},b_j^{(\ell)},\beta_j^{(\ell)},\xi_\ell,
        \sigma_{\alpha,\ell}^2:i,\ell,j\},\\
    \Theta_{\mathrm{mix}}
    &=\{c_q,\rho,\mu_l,\Sigma_l,\Delta_l,w_q,U_q,S_0:q,l\}.
\end{align*}
Using the hierarchical representation \eqref{eq:w-prior}--\eqref{eq:rho-prior},
the full joint posterior factors as
\begin{align}
&p(\Theta_{\mathrm{LSIRM}},\Theta_{\mathrm{mix}}\mid Y)
\notag\\
&\quad\propto
\prod_{\ell=1}^L\prod_{(i,j)\in\mathcal O_\ell}
    f_\ell\{y_{ij}^{(\ell)}\mid\eta_{ij}^{(\ell)},\xi_\ell\}
\notag\\
&\qquad\times
\prod_{q=1}^P
    N_d\!\bigl(z_q;\mu_{c_q}+\Delta_{c_q}U_q,\Sigma_{c_q}/w_q\bigr)
    \,\rho_{c_q}
    \,\mathrm{Gamma}(w_q;\nu_t/2,\nu_t/2)
    \,N_+(U_q;0,1/w_q)
\notag\\
&\qquad\times
\prod_{l=1}^{K^\star}
    p(\mu_l\mid\Sigma_l,S_0)\,p(\Sigma_l\mid S_0)\,
    N_d(\Delta_l;0,\tau_\Delta^2 I_d)
    \;p(\rho)\;p(S_0)
\notag\\
&\qquad\times
\prod_{i=1}^n N_d(a_i;0,I_d)
\prod_{\ell=1}^L\Bigl[\prod_{i=1}^n p(\alpha_i^{(\ell)}\mid\sigma_{\alpha,\ell}^2)\,
    p_\beta^{(\ell)}(\beta^{(\ell)})\Bigr]
\notag\\
&\qquad\times
\prod_{\ell=1}^L p(\sigma_{\alpha,\ell}^2)\,p(\xi_\ell),
\label{eq:full-posterior}
\end{align}
where $z_q=b_{j(q)}^{(\ell(q))}$ and $p_\beta^{(\ell)}$ is the
layer-specific item-difficulty prior factor (Gaussian for non-ordinal
layers; product of order-restricted Gaussians for ordinal layers).

\section{Residual augmentation}
\label{sec:residual}

The cornerstone of the SDB skew-$t$ sampler is the observation that
\emph{conditional on} $(\Delta_l,U_q)$, the model reduces to the
elliptical Student-$t$ mixture of the previous version on the residual
\begin{equation}
    r_q:=z_q-\Delta_{c_q}U_q.
    \label{eq:residual}
\end{equation}
Indeed, substituting $z_q=r_q+\Delta_{c_q}U_q$ in
\eqref{eq:z-given-latents} gives
\begin{equation}
    r_q\mid c_q=l,\mu_l,\Sigma_l,w_q
    \sim
    N_d\!\bigl(\mu_l,\;\Sigma_l/w_q\bigr),
    \label{eq:residual-Gaussian}
\end{equation}
which is the v11-style Gaussian scale mixture.  All NIW conjugate
machinery for $(\mu_l,\Sigma_l)$ -- weighted sufficient statistics,
posterior predictive density, collapsed allocation weights, NIW marginal
likelihood -- therefore carries over verbatim once $z_q$ is replaced by
$r_q$.

\section{NIW posterior on residuals}
\label{sec:niw-residual}

For any set $A\subseteq\{1,\dots,P\}$ of global item indices with current
residuals $\{r_q:q\in A\}$ and current scale weights $\{w_q:q\in A\}$,
define the weighted sufficient statistics
\begin{align}
    W_A &= \sum_{q\in A}w_q,
    \\
    \bar r_A^{\,w} &=
    \frac{1}{W_A}\sum_{q\in A}w_q\,r_q,
    \qquad W_A>0,
    \\
    S_A^{\,w} &= \sum_{q\in A}w_q\,
    (r_q-\bar r_A^{\,w})(r_q-\bar r_A^{\,w})\trans.
\end{align}
The weighted NIW posterior parameters for the items in $A$ are
\begin{align}
    \kappa_A &= \kappa_0+W_A,
    \label{eq:kappa-post}\\
    \nu_A &= \nu_0+|A|,
    \label{eq:nu-post}\\
    m_A &=
    \begin{cases}
    \dfrac{\kappa_0 m_0+W_A\bar r_A^{\,w}}{\kappa_0+W_A},
    & |A|>0,\\[1em]
    m_0, & |A|=0,
    \end{cases}
    \label{eq:m-post}\\
    S_A &=
    \begin{cases}
    S_0+S_A^{\,w}+
    \dfrac{\kappa_0 W_A}{\kappa_0+W_A}
    (\bar r_A^{\,w}-m_0)(\bar r_A^{\,w}-m_0)\trans,
    & |A|>0,\\[1em]
    S_0, & |A|=0.
    \end{cases}
    \label{eq:S-post}
\end{align}
The IW degrees-of-freedom update $\nu_A=\nu_0+|A|$ uses the item count
$|A|$ (not $W_A$) because each Gaussian factor in
\eqref{eq:residual-Gaussian} contributes one determinant
$|\Sigma_l|^{-1/2}$ regardless of $w_q$.

For component $l$, let $\mathcal J_l=\{q:c_q=l\}$ and $n_l=|\mathcal J_l|$.
Define
$(\kappa_l,\nu_l,m_l,S_l)
=(\kappa_{\mathcal J_l},\nu_{\mathcal J_l},m_{\mathcal J_l},S_{\mathcal J_l})$.
For leave-one-out summaries, $\mathcal J_{l,-q}=\{r:r\neq q,\,c_r=l\}$.

\section{NIW posterior predictive allocation update}
\label{sec:c-update}

Integrating out $(\mu_l,\Sigma_l)$ for the label allocation step uses the
NIW Student-$t$ posterior predictive on the residual $r_q$.  Because $r_q$
arises from \eqref{eq:residual-Gaussian} with covariance $\Sigma_l/w_q$,
the predictive scale is inflated by $w_q^{-1}$:
\begin{equation}
    r_q\mid r_{\mathcal J_{l,-q}},w_q
    \sim
    t_{\nu_{l,-q}-d+1}
    \!\left(
    m_{l,-q},\,
    \frac{\kappa_{l,-q}+w_q}{\kappa_{l,-q}\,w_q\,(\nu_{l,-q}-d+1)}\,S_{l,-q}
    \right).
    \label{eq:niw-predictive-skewt}
\end{equation}
Since $r_q=z_q-\Delta_l U_q$ when $q$ is tentatively assigned to cluster
$l$, the collapsed single-site label update is
\begin{equation}
    P(c_q=l\mid c_{-q},z,\Delta,U,w)
    \propto
    (n_{l,-q}+e_0)\;
    t_{\nu_{l,-q}-d+1}\!\Bigl(
        z_q-\Delta_l U_q;\,m_{l,-q},\,V_{l,-q}(w_q)
    \Bigr),
    \qquad l=1,\dots,K^\star,
    \label{eq:c-update-skewt}
\end{equation}
where $V_{l,-q}(w_q)=\dfrac{\kappa_{l,-q}+w_q}{\kappa_{l,-q}\,w_q\,(\nu_{l,-q}-d+1)}S_{l,-q}$
is the predictive scale matrix from \eqref{eq:niw-predictive-skewt}.

This is a Gibbs update on $c_q$, so the move is accepted with probability
one after sampling from the normalized weights \eqref{eq:c-update-skewt}.
Note that the candidate cluster label $l$ enters not only the NIW
sufficient statistics $(\kappa_{l,-q},\nu_{l,-q},m_{l,-q},S_{l,-q})$ but
also the residual location $z_q-\Delta_l U_q$, so $\Delta_l$ acts as a
cluster-specific shift in the predictive evaluation.

\section{Posterior draw of cluster parameters}
\label{sec:cluster-param-update}

After updating labels and the latents $(U,w)$, refresh the cluster
parameters in three sequential Gibbs steps.

\subsection{NIW draw of $(\mu_l,\Sigma_l)$}

Using the residual sufficient statistics of Section~\ref{sec:niw-residual}
computed at the \emph{current} $\Delta_l$ and $\{U_q:q\in\mathcal J_l\}$,
\begin{align}
    \Sigma_l\mid r,c
    &\sim\IW(\nu_l,S_l),
    \label{eq:Sigma-post}\\
    \mu_l\mid\Sigma_l,r,c
    &\sim N_d\!\left(m_l,\;\frac{1}{\kappa_l}\Sigma_l\right),
    \qquad l=1,\dots,K^\star.
    \label{eq:mu-post}
\end{align}
For empty components ($n_l=0$) this reduces to a draw from the NIW prior
\eqref{eq:Sigma-prior}--\eqref{eq:mu-prior}.

\subsection{Multivariate Gaussian Gibbs draw of $\Delta_l$}

Conditional on $(\mu_l,\Sigma_l)$ and on the per-item latents
$\{(U_q,w_q):q\in\mathcal J_l\}$, the contribution of cluster $l$ to the
likelihood of $\Delta_l$ from \eqref{eq:z-given-latents} is
\begin{equation}
    -\tfrac{1}{2}\sum_{q\in\mathcal J_l}w_q\,
    (z_q-\mu_l-\Delta_l U_q)\trans\Sigma_l^{-1}(z_q-\mu_l-\Delta_l U_q).
    \label{eq:Delta-likelihood}
\end{equation}
Expanding the quadratic in $\Delta_l$ and combining with the prior
\eqref{eq:Delta-prior} gives a $d$-dimensional Gaussian full conditional,
\begin{equation}
    \Delta_l\mid z,c,\mu,\Sigma,U,w
    \sim
    N_d\!\bigl(M_l,V_l\bigr),
    \label{eq:Delta-post}
\end{equation}
with precision and mean
\begin{align}
    V_l^{-1}
    &=
    \biggl(\sum_{q\in\mathcal J_l}w_q U_q^{\,2}\biggr)\,\Sigma_l^{-1}
    +\tau_\Delta^{-2}I_d,
    \label{eq:Delta-prec}\\
    M_l
    &=
    V_l\,\Sigma_l^{-1}\sum_{q\in\mathcal J_l}
    w_q U_q\,(z_q-\mu_l).
    \label{eq:Delta-mean}
\end{align}
For empty components, $V_l^{-1}=\tau_\Delta^{-2}I_d$ and $M_l=0$, so
$\Delta_l$ is drawn from the prior $N_d(0,\tau_\Delta^2 I_d)$.

\subsection{Half-normal Gibbs draw of $U_q$}

Conditional on cluster parameters and on $w_q$, the full conditional of
$U_q$ is obtained by combining the half-normal prior \eqref{eq:U-prior}
with the Gaussian likelihood \eqref{eq:z-given-latents}.  Both contribute
quadratics in $U_q$, giving a truncated normal:
\begin{equation}
    U_q\mid z,c,\mu,\Sigma,\Delta,w
    \sim
    N_+\!\bigl(\mu_{U,q},\,\sigma_{U,q}^{2}\bigr),
    \label{eq:U-post}
\end{equation}
where
\begin{align}
    \sigma_{U,q}^{-2}
    &=
    w_q\,\bigl(1+\Delta_{c_q}\trans\Sigma_{c_q}^{-1}\Delta_{c_q}\bigr),
    \label{eq:U-prec}\\
    \mu_{U,q}
    &=
    \sigma_{U,q}^{\,2}\,w_q\,
    \Delta_{c_q}\trans\Sigma_{c_q}^{-1}\,(z_q-\mu_{c_q}).
    \label{eq:U-mean}
\end{align}
The truncation to $(0,\infty)$ is sampled by inverse CDF: draw
$u\sim\mathrm{Uniform}\!\bigl(\Phi(-\mu_{U,q}/\sigma_{U,q}),1\bigr)$ and
return $\mu_{U,q}+\sigma_{U,q}\Phi^{-1}(u)$, where $\Phi$ is the standard
normal cdf.

\subsection{Gamma Gibbs draw of $w_q$}

Conditional on $(\Delta_{c_q},U_q,\mu_{c_q},\Sigma_{c_q})$, the full
conditional of $w_q$ from the Gamma prior \eqref{eq:w-prior} together with
the Gaussian likelihood \eqref{eq:z-given-latents} \emph{and} the
half-normal density $N_+(U_q;0,1/w_q)\propto w_q^{1/2}\exp(-w_q U_q^2/2)$
is again Gamma:
\begin{equation}
    w_q\mid z,c,\mu,\Sigma,\Delta,U
    \sim
    \mathrm{Gamma}\!\left(
    \frac{\nu_t+d+1}{2},\;
    \frac{\nu_t+D_q+U_q^{\,2}}{2}
    \right),
    \label{eq:w-post}
\end{equation}
where
\begin{equation}
    D_q
    =
    (z_q-\mu_{c_q}-\Delta_{c_q}U_q)\trans
    \Sigma_{c_q}^{-1}
    (z_q-\mu_{c_q}-\Delta_{c_q}U_q).
    \label{eq:Dq}
\end{equation}
The additional $+1$ in the shape and $+U_q^{\,2}/2$ in the rate (relative
to the elliptical Student-$t$ version) come exactly from the half-normal
contribution.

\subsection{Wishart Gibbs draw of $S_0$}

After all $\Sigma_l$ have been refreshed, draw the cluster-scale
hyperparameter from \eqref{eq:S0-update}:
\begin{equation}
    S_0\mid\{\Sigma_l\}
    \sim
    W\!\left(
    \nu_{S_0}+K^\star\nu_0,\;
    \bigl(\Lambda_{S_0}^{-1}+\textstyle\sum_l\Sigma_l^{-1}\bigr)^{-1}
    \right).
    \label{eq:S0-gibbs}
\end{equation}
The cached log-determinant $\log|S_0|$, used in the NIW marginal likelihood
\eqref{eq:niw-marginal} for split--merge moves, must be refreshed after
every $S_0$ draw.

\subsection{Mixture weight $\rho$ (storage only)}

For monitoring,
\begin{equation}
    \rho\mid c
    \sim
    \Dir(e_0+n_1,\dots,e_0+n_{K^\star}).
    \label{eq:rho-post}
\end{equation}
The draw of $\rho$ is not needed for the collapsed label update
\eqref{eq:c-update-skewt}.

\section{Collapsed partition target for split--merge moves}
\label{sec:collapsed-partition}

Integrating out $\rho$, $(\mu_l,\Sigma_l)$, $w_q$, and $U_q$ would give a
partition target on $c$ alone, but the integrals over the per-item latents
$(w_q,U_q)$ do not have a closed form for SDB skew-$t$.  We instead
condition on $(\Delta_l,U_q,w_q)$ and integrate only $\rho$ and
$(\mu_l,\Sigma_l)$.  The resulting target on $c$ is
\begin{equation}
    \pi(c\mid r,\Delta)
    \propto
    \prod_{l=1}^{K^\star}
    \frac{\Gamma(n_l+e_0)}{\Gamma(e_0)}
    \,m(r_{\mathcal J_l}),
    \label{eq:partition-target}
\end{equation}
where the NIW marginal likelihood of the residuals in cluster $l$ is
\begin{equation}
    m(r_A)
    =
    \int\!\!\int
    \prod_{q\in A}N_d(r_q;\mu,\Sigma/w_q)
    \,p(\mu,\Sigma)\,d\mu\,d\Sigma,
    \label{eq:niw-marginal-residual}
\end{equation}
with $(\mu,\Sigma)\sim\mathrm{NIW}(m_0,\kappa_0,\nu_0,S_0)$.  The log
marginal in closed form is
\begin{align}
    \log m(r_A)
    &=-\frac{n_A d}{2}\log\pi
    +\frac{d}{2}\bigl\{\log\kappa_0-\log\kappa_A\bigr\}
    +\frac{\nu_0}{2}\log|S_0|
    -\frac{\nu_A}{2}\log|S_A|
    \notag\\
    &\qquad
    +\log\Gamma_d\!\left(\frac{\nu_A}{2}\right)
    -\log\Gamma_d\!\left(\frac{\nu_0}{2}\right)
    +\frac{d}{2}\sum_{q\in A}\log w_q,
    \label{eq:niw-marginal}
\end{align}
where
\begin{equation}
    \log\Gamma_d(a)
    =
    \frac{d(d-1)}{4}\log\pi
    +\sum_{h=1}^d\log\Gamma\!\left(a+\frac{1-h}{2}\right)
\end{equation}
is the multivariate Gamma function and $n_A=|A|$.  The
$\frac{d}{2}\sum_q\log w_q$ term collects the Jacobian factors from each
$w_q$-scaled Gaussian factor in \eqref{eq:niw-marginal-residual} and is
\emph{constant in any partition move} that preserves the item set, so it
cancels in the split/merge Metropolis--Hastings ratio.

\section{Metropolis--Hastings updates for LSIRM parameters}
\label{sec:lsirm-mh}

All LSIRM parameters with non-conjugate full conditionals are updated by
random-walk Metropolis--Hastings.  All proposals below are symmetric
Gaussian random walks with item- or parameter-specific scale.

\subsection{Respondent effect $\alpha_i^{(\ell)}$}

Propose
$\alpha_i^{(\ell)\prime}=\alpha_i^{(\ell)}+\varepsilon_\alpha$,
$\varepsilon_\alpha\sim N(0,s_\alpha^2)$.  Let
$\eta_{ij}^{(\ell)\prime}$ be the linear predictor after the proposal.
The log acceptance ratio is
\begin{align}
    \log R_\alpha
    &=
    \sum_{j:(i,j)\in\mathcal O_\ell}
    \!\bigl[\log f_\ell\{y_{ij}^{(\ell)}\mid\eta_{ij}^{(\ell)\prime},\xi_\ell\}
        -\log f_\ell\{y_{ij}^{(\ell)}\mid\eta_{ij}^{(\ell)},\xi_\ell\}\bigr]
    \notag\\
    &\qquad
    -\frac{1}{2\sigma_{\alpha,\ell}^{\,2}}
    \bigl\{(\alpha_i^{(\ell)\prime})^2-(\alpha_i^{(\ell)})^2\bigr\}.
    \label{eq:alpha-mh}
\end{align}

\subsection{Item difficulty $\beta_j^{(\ell)}$}

For non-ordinal layers, propose
$\beta_j^{(\ell)\prime}=\beta_j^{(\ell)}+\varepsilon_\beta$,
$\varepsilon_\beta\sim N(0,s_\beta^2)$.  With the fixed-scale prior
\eqref{eq:beta-prior},
\begin{align}
    \log R_\beta
    &=
    \sum_{i:(i,j)\in\mathcal O_\ell}
    \!\bigl[\log f_\ell\{y_{ij}^{(\ell)}\mid\eta_{ij}^{(\ell)\prime},\xi_\ell\}
        -\log f_\ell\{y_{ij}^{(\ell)}\mid\eta_{ij}^{(\ell)},\xi_\ell\}\bigr]
    \notag\\
    &\qquad
    -\frac{1}{2\sigma_{\beta,\ell}^{\,2}}
    \bigl\{(\beta_j^{(\ell)\prime})^2-(\beta_j^{(\ell)})^2\bigr\}.
    \label{eq:beta-mh}
\end{align}
For ordinal layers, propose each threshold $\beta_{j,k}^{(\ell)\prime}$
by Gaussian random walk; reject the proposal outright if the order
constraint $\beta_{j,1}^{(\ell)}<\cdots<\beta_{j,K_\ell-1}^{(\ell)}$ is
violated; otherwise compute the analogue of \eqref{eq:beta-mh} with the
layer-$\ell$ ordinal likelihood
\eqref{eq:lik-ord-cum}--\eqref{eq:lik-ord}.

\subsection{Distance weight $\gamma_\ell$}

Propose $\log\gamma_\ell'=\log\gamma_\ell+\varepsilon_\gamma$,
$\varepsilon_\gamma\sim N(0,s_{\log\gamma}^2)$.  The log acceptance ratio
is the log-likelihood difference of all observations in layer $\ell$ at
$\gamma_\ell'$ versus $\gamma_\ell$ plus the log-normal prior difference
on $\log\gamma_\ell$.

\subsection{Count-layer overdispersion $\kappa_j$}

Propose $\log\kappa_j'=\log\kappa_j+\varepsilon_\kappa$,
$\varepsilon_\kappa\sim N(0,s_{\log\kappa}^2)$.  The log acceptance ratio
is the sum of negative-binomial log-likelihood differences over
respondents and the log-normal prior difference \eqref{eq:kappa-prior}.

\subsection{Respondent latent position $a_i$}

Propose $a_i'=a_i+\varepsilon_a$, $\varepsilon_a\sim N_d(0,s_a^2 I_d)$,
and aggregate likelihood contributions over every layer:
\begin{align}
    \log R_a
    &=
    \sum_{\ell=1}^L\sum_{j:(i,j)\in\mathcal O_\ell}
    \!\bigl[\log f_\ell\{y_{ij}^{(\ell)}\mid\eta_{ij}^{(\ell)\prime},\xi_\ell\}
        -\log f_\ell\{y_{ij}^{(\ell)}\mid\eta_{ij}^{(\ell)},\xi_\ell\}\bigr]
    \notag\\
    &\qquad
    -\tfrac{1}{2}\bigl\{\norm{a_i'}_2^2-\norm{a_i}_2^2\bigr\},
    \label{eq:a-mh}
\end{align}
where the prior variance is fixed at $1$.

\subsection{Item latent position $b_j^{(\ell)}$}
\label{subsec:b-mh}

Let $q=q(\ell,j)$, with current cluster label $c_q$ and current latents
$(U_q,w_q)$.  Propose $b_j^{(\ell)\prime}=b_j^{(\ell)}+\varepsilon_b$,
$\varepsilon_b\sim N_d(0,s_b^2 I_d)$.  The log acceptance ratio combines
the LSIRM likelihood difference and the conditional Gaussian prior from
\eqref{eq:z-given-latents}:
\begin{align}
    \log R_b
    &=
    \sum_{i:(i,j)\in\mathcal O_\ell}
    \!\bigl[\log f_\ell\{y_{ij}^{(\ell)}\mid\eta_{ij}^{(\ell)\prime},\xi_\ell\}
        -\log f_\ell\{y_{ij}^{(\ell)}\mid\eta_{ij}^{(\ell)},\xi_\ell\}\bigr]
    \notag\\
    &\qquad
    -\tfrac{1}{2}w_q\,
    \bigl\{
    (b_j^{(\ell)\prime}-\mu_{c_q}-\Delta_{c_q}U_q)\trans
    \Sigma_{c_q}^{-1}
    (b_j^{(\ell)\prime}-\mu_{c_q}-\Delta_{c_q}U_q)
    \notag\\
    &\hspace{4em}
    -(b_j^{(\ell)}-\mu_{c_q}-\Delta_{c_q}U_q)\trans
    \Sigma_{c_q}^{-1}
    (b_j^{(\ell)}-\mu_{c_q}-\Delta_{c_q}U_q)
    \bigr\}.
    \label{eq:b-mh-skewt}
\end{align}
This is the conditional-on-$(c_q,\mu,\Sigma,\Delta,U,w)$ MH ratio.  The
cluster assignment $c_q$ enters \eqref{eq:b-mh-skewt} through the
anchor $\mu_{c_q}+\Delta_{c_q}U_q$; the cluster scatter $\Sigma_{c_q}$
controls the prior strength; and the scale latent $w_q$ inflates or
deflates that strength inversely.

\section{Conjugate update for the respondent-effect variance}

The only random-effect variance that is sampled is the layer-specific
respondent-effect variance $\sigma_{\alpha,\ell}^2$, with full conditional
\begin{equation}
    \sigma_{\alpha,\ell}^2\mid\alpha^{(\ell)}
    \sim
    \IG\!\left(
    A_{\alpha,\ell}+\frac{n}{2},\;
    B_{\alpha,\ell}+\frac{1}{2}\sum_{i=1}^n(\alpha_i^{(\ell)})^2
    \right).
    \label{eq:sigma-alpha-update}
\end{equation}
The respondent-position prior variance is fixed at $1$ \eqref{eq:a-prior};
the item-difficulty scales $\sigma_{\beta,\ell}^{\,2}$ are fixed; the
continuous-layer error scale $\sigma_0^2$ has its own inverse-gamma full
conditional (standard Gaussian-likelihood conjugate update).

\section{Jain--Neal restricted Gibbs split--merge moves}
\label{sec:split-merge}

Split--merge moves are Metropolis--Hastings updates on $c$ conditional on
the current residuals $r=(r_1,\dots,r_P)$ (Section~\ref{sec:residual}) and
on the current $\Delta_l$ values.  The target is the collapsed partition
density in \eqref{eq:partition-target}.  All residual evaluations use the
candidate cluster's $\Delta_l$ (so that an item proposed for cluster $l$
contributes residual $z_q-\Delta_l U_q$).

\subsection{Anchor selection}

Select two distinct global item indices $u$ and $v$ uniformly from
$\{1,\dots,P\}$.  If $c_u=c_v$, attempt a \emph{split}; otherwise attempt
a \emph{merge}.  The anchor-selection probability is symmetric and cancels
in the MH ratio.

\subsection{Split proposal}

Suppose $c_u=c_v=a$ and let $A=\{q:c_q=a\}$.  Choose an empty component
label $b\in\{l:n_l=0\}$ uniformly at random; if no empty label exists,
skip the split.  The newly opened component inherits its parent's skew
direction by setting $\Delta_b^{\mathrm{prop}}=\Delta_a$.

Initialize two temporary clusters by forcing $c_u'=a$, $c_v'=b$.  For
each $q\in A\setminus\{u,v\}$, assign $q$ to either $a$ or $b$ by a
restricted residual-collapsed Gibbs scan with weights
\begin{align}
    w_q(a)&\propto
    (n_{a,-q}^{\mathrm{temp}}+e_0)\,
    t_{\nu_{a,-q}-d+1}\!\bigl(
        z_q-\Delta_a U_q;\,m_{a,-q}^{\mathrm{temp}},V_{a,-q}^{\mathrm{temp}}(w_q)
    \bigr),
    \label{eq:split-w-a}\\
    w_q(b)&\propto
    (n_{b,-q}^{\mathrm{temp}}+e_0)\,
    t_{\nu_{b,-q}-d+1}\!\bigl(
        z_q-\Delta_b^{\mathrm{prop}} U_q;\,m_{b,-q}^{\mathrm{temp}},V_{b,-q}^{\mathrm{temp}}(w_q)
    \bigr),
    \label{eq:split-w-b}
\end{align}
where the NIW posterior summaries
$(\kappa,\nu,m,S)^{\mathrm{temp}}_{a,-q}$,
$(\kappa,\nu,m,S)^{\mathrm{temp}}_{b,-q}$ are computed from the
\emph{current} temporary assignments of items into clusters $a$ and $b$,
respectively, using residuals
$\{z_{q'}-\Delta_a U_{q'}\}$ and $\{z_{q'}-\Delta_b^{\mathrm{prop}}U_{q'}\}$.
Let $p_q^{\mathrm{fwd}}$ be the normalized probability of the assignment
actually selected.  Then
\begin{equation}
    \log q(c'\mid c)
    =
    -\log E(c)
    +\sum_{q\in A\setminus\{u,v\}}\log p_q^{\mathrm{fwd}},
    \label{eq:split-fwd}
\end{equation}
where $E(c)$ is the number of empty components before the split.  The
reverse move is a deterministic merge, so $\log q(c\mid c')=0$.  The
split acceptance probability is
\begin{equation}
    \alpha_{\mathrm{split}}
    =
    \min\!\left\{1,
    \exp\bigl[
    \log\pi(c'\mid r,\Delta)-\log\pi(c\mid r,\Delta)
    -\log q(c'\mid c)
    \bigr]\right\}.
    \label{eq:split-acc}
\end{equation}
If the move is accepted, the next sweep's Gibbs step for $\Delta_b$
\eqref{eq:Delta-post} re-estimates the new cluster's skew direction from
its newly assigned items.

\subsection{Merge proposal}

Suppose $c_u=a$ and $c_v=b$, $a\neq b$.  Let $A=\{q:c_q=a\}$ and
$B=\{q:c_q=b\}$.  The proposal merges $B$ into $A$:
$c_q'=a$ for all $q\in A\cup B$; the merged cluster keeps $\Delta_a$.
Thus $\log q(c'\mid c)=0$.

The reverse proposal would split $A\cup B$ back into $A$ and $B$ using
$u,v$ as anchors with restricted-Gibbs scan probabilities analogous to
\eqref{eq:split-w-a}--\eqref{eq:split-w-b}, with $b$'s temporary skew
direction set to the merged-state $\Delta_a$ (so the reverse path is
well-defined).  Let $p_q^{\mathrm{rev}}$ be the probability of assigning
$q$ to its original side during this reverse scan.  Then
\begin{equation}
    \log q(c\mid c')
    =
    -\log E(c')
    +\sum_{q\in(A\cup B)\setminus\{u,v\}}\log p_q^{\mathrm{rev}},
    \label{eq:merge-rev}
\end{equation}
where $E(c')$ is the number of empty components after the merge.  The
merge acceptance probability is
\begin{equation}
    \alpha_{\mathrm{merge}}
    =
    \min\!\left\{1,
    \exp\bigl[
    \log\pi(c'\mid r,\Delta)-\log\pi(c\mid r,\Delta)
    +\log q(c\mid c')
    \bigr]\right\}.
    \label{eq:merge-acc}
\end{equation}

\paragraph{Treatment of $\Delta_l$ in split--merge.}
Both the forward split and the reverse split (inside the merge MH ratio)
condition on the current $\Delta_l$ values rather than marginalizing them.
In split, the new component inherits its parent's $\Delta$; in merge, the
surviving component keeps its $\Delta$; the Gibbs step
\eqref{eq:Delta-post} subsequently refreshes both.  This conditional-on-
$\Delta$ scheme preserves the target $\pi(c\mid r,\Delta)$ of
\eqref{eq:partition-target} as the partition-only stationary distribution
of the MH move; the full joint posterior \eqref{eq:full-posterior} is
preserved because every conditioning variable is refreshed elsewhere in
the sweep.

\section{Posterior cluster summaries}

Because labels are not identifiable, cluster summaries should be
label-switching invariant.  The primary summary is the posterior
similarity matrix
\begin{equation}
    C_{qr}=P(c_q=c_r\mid Y),
    \qquad q,r=1,\dots,P,
\end{equation}
updated online by
$C_{qr}^{\mathrm{raw}}\leftarrow C_{qr}^{\mathrm{raw}}+\ind(c_q=c_r)$
after burn-in.  The trace of the occupied-component count
$K_+^{(t)}=|\{l:n_l^{(t)}>0\}|$ is monitored for chain mixing.
Point-estimate partitions are obtained by minimising the Binder loss or
the variation-of-information (VI) loss over the posterior similarity
matrix.

The skew directions $\{\Delta_l\}$ admit additional summaries, but they
are subject to the same rotation ambiguity as $(\mu_l,\Sigma_l)$
\eqref{eq:rigid-motion}, so visualization of $\Delta_l$ should be paired
with the post-hoc Procrustes alignment of
Section~\ref{sec:rotation-invariance}.

\section{Sweep order and validity}
\label{sec:sweep}

A single MCMC sweep updates each block of parameters once.  The order is
\begin{enumerate}[label=(\arabic*),leftmargin=2em]
    \item Layer-specific nuisance latents: continuous-layer
    $\{\lambda_{ij}^{(\ell)}\}$ Gibbs;
    \item Respondent effects $\{\alpha_i^{(\ell)}\}$ MH per \eqref{eq:alpha-mh};
    \item Respondent positions $\{a_i\}$ MH per \eqref{eq:a-mh};
    \item Item difficulties $\{\beta_j^{(\ell)}\}$ MH per \eqref{eq:beta-mh};
    \item Item positions $\{b_j^{(\ell)}\}$ MH per \eqref{eq:b-mh-skewt};
    \item Ordinal thresholds (when applicable), distance weights
    $\{\gamma_\ell\}$, count overdispersions $\{\kappa_j\}$ MH;
    \item Continuous-layer scale $\sigma_0^2$ inverse-gamma Gibbs;
    respondent-effect variances $\sigma_{\alpha,\ell}^{2}$ inverse-gamma
    Gibbs per \eqref{eq:sigma-alpha-update};
    \item Allocation Gibbs $\{c_q\}$ per \eqref{eq:c-update-skewt};
    \item Jain--Neal split--merge proposals
    \eqref{eq:split-acc}, \eqref{eq:merge-acc}, repeated $M_{\mathrm{SM}}$
    times;
    \item NIW Gibbs $\{(\mu_l,\Sigma_l)\}$ per
    \eqref{eq:Sigma-post}--\eqref{eq:mu-post};
    \item Skew direction Gibbs $\{\Delta_l\}$ per \eqref{eq:Delta-post};
    \item Half-normal Gibbs $\{U_q\}$ per \eqref{eq:U-post};
    \item Gamma Gibbs $\{w_q\}$ per \eqref{eq:w-post};
    \item Wishart Gibbs $S_0$ per \eqref{eq:S0-gibbs};
    \item (Monitoring only) Dirichlet draw of $\rho$ per \eqref{eq:rho-post}.
\end{enumerate}

\paragraph{Validity.}
Each block is a proper full-conditional Gibbs or a Metropolis--Hastings
move targeting the conditional distribution implied by the joint posterior
\eqref{eq:full-posterior}.  In particular, the partition move
$\pi(c\mid r,\Delta)$ targets the marginal distribution of $c$ after
integrating $\rho$, $\mu$, $\Sigma$ out of the joint, conditional on the
\emph{current} $\Delta$ and on the residuals $r=z-\Delta U$ (which depend
on the current $U$); since $\Delta$, $U$, and $w$ are refreshed in
subsequent blocks before they are reused, the sweep order
(1)--(15) is a valid Gibbs scan of the joint posterior.

\section{Hyperparameter setting criteria}
\label{sec:hyperparams}

The hyperparameters governing the skew-$t$ mixture and its priors are set
as follows.  All scales below refer to the implementation's default
latent-space scale (item-position norms typically in $[0,3]$ after
Procrustes alignment to the posterior-mean configuration).

\paragraph{Student-$t$ degrees of freedom $\nu_t$ (fixed).}
$\nu_t>2$ ensures finite within-cluster variance.  Heavier tails (smaller
$\nu_t$) tolerate boundary items but reduce the cluster center's
identifying power.  Recommended range $\nu_t\in[3,10]$; default
$\nu_t=3$.  This is treated as fixed rather than estimated because
$\nu_t$ and the cluster-shape concentration $\nu_0$ are weakly
identified jointly.

\paragraph{Skew direction prior scale $\tau_\Delta$ (fixed).}
$\tau_\Delta$ controls the prior strength on cluster asymmetry.
$\tau_\Delta\to 0$ collapses the model to elliptical Student-$t$;
$\tau_\Delta$ large allows arbitrarily strong skew.  Recommended range
$\tau_\Delta\in[0.5,2.0]$; default $\tau_\Delta=1.0$, matching the
expected within-cluster spread in the latent space.  For datasets where
prior knowledge suggests symmetric clusters, $\tau_\Delta=0.1$ provides
shrinkage toward $\Delta_l=0$ while allowing data to override if the
posterior signal is strong.

\paragraph{NIW concentration on $\mu_l$: $\kappa_0$ (fixed).}
$\kappa_0$ is the prior effective sample size for the cluster center.
Small $\kappa_0$ (e.g.\ $10^{-3}$) gives essentially no prior pull on
$\mu_l$, which can let empty-component $\mu_l$ draws explode under the
$N(m_0,\Sigma_l/\kappa_0)$ prior.  Recommended $\kappa_0\in[0.1,2]$;
default $\kappa_0=1.0$, which gives a cluster-center prior standard
deviation comparable to the expected within-cluster spread.

\paragraph{NIW concentration on $\Sigma_l$: $\nu_0$ (fixed).}
$\nu_0>d-1$ ensures the IW prior is proper; the prior mean is
$E[\Sigma_l]=S_0/(\nu_0-d-1)$.  Larger $\nu_0$ pins $\Sigma_l$ closer to
this mean; smaller $\nu_0$ allows more data-driven cluster shape.
Recommended $\nu_0=d+10$; the default $\nu_0=12$ (at $d=2$) is moderately
informative without overconstraining cluster anisotropy.

\paragraph{NIW prior mean $m_0$ (fixed).}
$m_0=0$ is the canonical choice when Procrustes alignment is used to
anchor the latent space at the origin.  No data-driven $m_0$ is needed.

\paragraph{NIW prior scale $S_0$ (initial value, then Wishart-sampled).}
The initial $S_0^{(0)}=s_0 I_d$ sets the prior expected cluster
covariance $E[\Sigma_l]=S_0^{(0)}/(\nu_0-d-1)$.  Recommended $s_0=1$
(again on the latent-space scale).  The Wishart hyperprior
\eqref{eq:S0-hyperprior} with $\nu_{S_0}=d+2$ and
$\Lambda_{S_0}=S_0^{(0)}/\nu_{S_0}$ gives $E[S_0]=S_0^{(0)}$ while
allowing data-driven adjustment.  Setting $\nu_{S_0}\to\infty$ holds
$S_0$ fixed.

\paragraph{Dirichlet concentration $e_0$ (fixed).}
The symmetric Dirichlet \eqref{eq:rho-prior} with small $e_0$
(``overfitted finite mixture'' regime of Rousseau and Mengersen, 2011)
asymptotically empties redundant components when $K^\star$ exceeds the
true cluster count.  The Rousseau--Mengersen safe range is
$e_0<d_l/2$ where $d_l=d+d(d+1)/2$ is the per-component parameter
count of the Gaussian core (ignoring $\Delta_l$); for $d=2$ this gives
$e_0<2.5$.  Recommended $e_0\in[0.05,0.5]$; default $e_0=0.1$.

\paragraph{Truncation level $K^\star$ (fixed).}
$K^\star$ is the maximum number of mixture components.  The chain
automatically empties components beyond the true cluster count when
$e_0$ is in the overfitted regime, so $K^\star$ can be set generously.
A practical choice is $K^\star=\min(P,\,K_{\mathrm{expected}}+3)$, with
$K_{\mathrm{expected}}$ a conservative upper bound on the number of
substantively meaningful clusters.

\paragraph{Number of split--merge attempts $M_{\mathrm{SM}}$ per sweep.}
Each restricted-Gibbs split--merge proposal has lower per-attempt
acceptance under the joint LSIRM coupling than a stand-alone clustering
model because the residuals $r_q$ are updated only by the slower
random-walk MH on $b_j^{(\ell)}$ rather than by direct Gibbs.  Increasing
$M_{\mathrm{SM}}$ compensates: $M_{\mathrm{SM}}=100$ typically gives a
$\geq 5\times$ improvement in the effective number of accepted moves per
sweep relative to $M_{\mathrm{SM}}=5$.

\paragraph{MCMC length, burn-in, and thinning.}
The MH-heavy LSIRM block requires longer chains than a pure Gibbs
clustering model.  Recommended initial settings are $n_{\mathrm{iter}}\in
[30000,100000]$, burn-in $\in[10000,30000]$, thin $\in[5,10]$, with the
chain length doubled if posterior-similarity-matrix convergence diagnostics
fail.  The burn-in for the FMC block can be shorter (e.g.\
$0.25\times\mathrm{burnin}_{\mathrm{LSIRM}}$) because $\{c_q\}$ converges
faster than the latent positions.

\paragraph{Random-walk MH proposal scales.}
The proposal scales $(s_a,s_b,s_\alpha,s_\beta,s_{\log\gamma},
s_{\log\kappa})$ should be tuned to per-parameter MH acceptance rates of
$0.30$--$0.40$ during a preliminary run.  In the implementation, scales
that yield acceptance outside $[0.15,0.55]$ are halved (if too low) or
doubled (if too high).

\section*{References}
\addcontentsline{toc}{section}{References}

\begin{itemize}[leftmargin=1.5em]
    \item Azzalini, A., and Dalla Valle, A.\ (1996).
    The multivariate skew-normal distribution.
    \emph{Biometrika}, 83(4), 715--726.

    \item De~Carolis, G., Kang, S., and Jeon, M.\ (2025).
    Latent space graded response model (LSGRM) for ordinal item response
    data.  Working paper.

    \item Gower, J.~C.\ (1975).  Generalized procrustes analysis.
    \emph{Psychometrika}, 40(1), 33--51.

    \item Jain, S., and Neal, R.~M.\ (2004).
    A split--merge Markov chain Monte Carlo procedure for the
    Dirichlet process mixture model.
    \emph{Journal of Computational and Graphical Statistics}, 13(1),
    158--182.

    \item Lange, K.~L., Little, R.~J.~A., and Taylor, J.~M.~G.\ (1989).
    Robust statistical modeling using the $t$ distribution.
    \emph{Journal of the American Statistical Association}, 84(408),
    881--896.

    \item Murphy, K.~P.\ (2007).
    Conjugate Bayesian analysis of the Gaussian distribution.
    Technical report, University of British Columbia.

    \item Rousseau, J., and Mengersen, K.\ (2011).
    Asymptotic behaviour of the posterior distribution in overfitted
    mixture models.
    \emph{Journal of the Royal Statistical Society, Series B}, 73(5),
    689--710.

    \item Sahu, S.~K., Dey, D.~K., and Branco, M.~D.\ (2003).
    A new class of multivariate skew distributions with applications to
    Bayesian regression models.
    \emph{The Canadian Journal of Statistics}, 31(2), 129--150.

    \item Sch\"onemann, P.~H.\ (1966).
    A generalized solution of the orthogonal procrustes problem.
    \emph{Psychometrika}, 31(1), 1--10.

    \item van~Dyk, D.~A., and Park, T.\ (2008).
    Partially collapsed Gibbs samplers: theory and methods.
    \emph{Journal of the American Statistical Association}, 103(482),
    790--796.
\end{itemize}

\end{document}
